\section{Audit of the existing implementation}
\label{sec:audit}

\subsection{Target identified, and why}

\texttt{quantdata/dysanos/} contains several largely independent DYSANOS
subsystems: \texttt{flex/}, \texttt{grid/gendysanos}, \texttt{grid/rndysanos},
\texttt{implied/}, \texttt{localvolbysanos/}, \texttt{tangent\_levy/}. This
report targets \texttt{flex/}, the only subsystem whose language (``AR(1)'',
weekly, joint DLV+grid+smoothing+dynamics estimation) matches the assignment's
description exactly, and the only one with an explicit, already-present
pathwise HJM/roll-residual diagnostic directly on point for the pricing
identity (Sec.~\ref{sec:sec23}). No trained checkpoint exists in this
repository checkout; numbers below fit fresh against local fixture/synthetic
data, stated as such throughout.

\subsection{Latent state and decoder}

\texttt{bound\_native\_raw} maps an unconstrained raw torch tensor to
\texttt{NativeParameters(dlvsqs, forward\_vols, eta)}: \texttt{dlvsqs}
$(T{\times}N)$ bounded to $[0.01^2,3.0^2]$, \texttt{forward\_vols} $(T,)$
bounded to $[0.01,1.0]$, \texttt{eta} (fixed or bounded $[0.01,0.50]$) via
sigmoid/interval transforms. \texttt{FlexSanosDecoder} then runs the
Wissel/Magnus finite-grid recursion directly on \texttt{dlvsqs} (confirmed
from the weight construction \texttt{w\_plus}/\texttt{w\_minus}, not assumed
from naming).

\subsection{Two term structures, and the variance/std convention}

\textbf{Grid geometry} (\texttt{num\_inner\_strikes}, \texttt{nms\_range},
\texttt{num\_model\_expiries}, the calendar roll schedule) is fixed per run.
\textbf{Smoothing} (\texttt{eta}, \texttt{forward\_vols}) is, materially,
\emph{part of the fitted state} -- \texttt{forward\_vols} is jointly
estimated, not a fixed side channel. Traced directly through
\texttt{price\_first\_expiry\_calls\_from\_parameters} $\to$
\texttt{\_black\_call}: the kernel dispersion is
\texttt{forward\_vol*eta*sqrt(dt)}, passed to the argument literally named
\texttt{sqrt\_variance} and used directly (not squared) in $d_1$ -- a standard
deviation, with \texttt{eta} itself a \emph{dimensionless multiplier} on
\texttt{forward\_vol}, not a standard deviation on its own. This differs
materially from other SANOS conventions in this research programme (where
$\eta$ \emph{is} the smoothing standard deviation directly) and must be
preserved exactly when porting any restriction between them.

\subsection{Statistical transition, front-expiry law, roll conventions}

\texttt{weekly\_dynamics.py}: a PCA-reduced, per-expiry-block, hierarchical
VAR(1) (\texttt{HierarchicalBlockFit(intercept, own\_lag, previous\_lag)}) at
a fixed 5-business-day stride, ridge-regularised, with a joint innovation
covariance and a spot process that loads on the surface innovation via
\texttt{weekly\_spot\_front\_beta} plus an orthogonal residual. A
continuous-time generator $K=-\log(A)/\Delta t$ is computed and \emph{gated}
(condition-number ceiling, exponential round-trip check) -- a P-side
stability diagnostic only; $\Q$ does not appear anywhere in this file.
\texttt{fixed\_expiry\_ar1.py}/\texttt{last\_expiry\_ar1.py} handle the DLV
tower's calendar-exact roll; the front-expiry settlement payoff used in the
codebase's own audit is the ordinary intrinsic value
(\texttt{relu(R-k)}), not a grid-pinned substitute.

\subsection{Existing pricing-drift infrastructure (the central prior-work finding)}

\texttt{WeeklySanosSimulator.hjm\_roll\_residuals} already computes, pathwise,
exactly the discrete martingale/roll residual Sec.~\ref{sec:sec23} formalises
and tests: \texttt{R*c\_next[m-1](k/R) - c\_current[m](k)} for non-front
expiries, with the correct terminal payoff at the front. Its only caller,
\texttt{weekly\_measure\_training.py}, already implements a sequential-Monte-
Carlo/gradient-based fit of a ``target-measure'' kernel
(\texttt{NormalWeeklyMeasureKernel}/\texttt{GenericWeeklyMeasureKernel}) via
\texttt{hjm\_roll\_moment\_loss} with instrumented conditioning -- a
substantial prior attempt at $\Pp\to\Q$ construction, used directly in
Sec.~\ref{sec:sec23} rather than reinvented. A family of
\texttt{run\_*\_q\_existence.py} scripts investigates the same admissibility
question via Wasserstein-2 roll conditions against real historical fits; their
result artefacts live in a training cache not present in this checkout and
were not read here -- flagged explicitly rather than silently worked around.

\textbf{The clearest gap}: no price-preserving reset (an exact constrained
inverse under admissibility constraints, CRC step E) exists anywhere in
\texttt{flex/}. This is what Sec.~\ref{sec:sec5} builds.

\subsection{Reproducible baseline}

Python~3.12, fresh \texttt{.venv312} built from \texttt{pyproject.toml},
\texttt{QUANTDATA\_TEST\_MODE=1}: \texttt{pytest
tests/test\_dysanos\_torch\_simulator.py tests/test\_dysanos\_market.py} --
\textbf{18/18 pass}. No file outside \texttt{flex/crc\_extension/} was
modified.
